% ─────────────────────────── COVER PAGE ───────────────────────────
\begin{titlepage}
\centering
\vspace{-0.8cm}
\includegraphics[width=3.2cm]{aast.png}\par
\vspace{0.4cm}
{\large\bfseries Arab Academy for Science, Technology and Maritime Transport\par}
\vspace{0.4cm}
{\normalsize College of Engineering and Technology\par}
{\normalsize Computer Engineering Department\par}
\vspace{1.2cm}
{\normalsize B.Sc.\ Final Year Project\par}
\vspace{1.8cm}
{\LARGE\bfseries \projecttitle\par}
\vspace{2.2cm}
{\normalsize\bfseries Presented by:\par}
\vspace{0.4cm}
\begin{tabular}{ll}
Omar Hossam Eldin Mohamed Hafez & Reg.\ No.\ 20101737\\
Mohamed Sameh Mohamed Mohamed El Bassuni & Reg.\ No.\ 18100896\\
Mohamed Mohamed Talaat Abokhatwa & Reg.\ No.\ 12103172\\
Ahmed Abdelaal Saad & Reg.\ No.\ 16101027\\
Adam Mohamed Ibrahim & Reg.\ No.\ 20101978\\
\end{tabular}
\vspace{1.8cm}

{\normalsize\bfseries Supervised by:\par}
\vspace{0.4cm}
{\normalsize Yahia Elderiny\par}
{\normalsize Rania Kadry\par}
\vfill
{\normalsize July --- 2026\par}
\end{titlepage}

\pagenumbering{roman}
\setcounter{page}{2}

% ─────────────────────────── DECLARATION ───────────────────────────
\clearpage
\section*{Declaration}
\addcontentsline{toc}{section}{Declaration}

I hereby certify that this report, which I am now submitting for evaluation on
the study program leading to the award of Bachelor of Science in Computer
Engineering, is all my own work and contains no plagiarism. I comply with the
following terms when submitting this report:

Any text, diagrams or other material copied from other sources (including, but
not limited to, books, journals, and the internet) have been clearly
acknowledged and cited, followed by the reference number used, either in the
text or in a footnote/endnote. The details of the used references that are
listed at the end of the report conform to the referencing style dictated by
the final year project template and are, to my knowledge, accurate and
complete.

I have read the sections on referencing and plagiarism in the final year
project template. I understand that plagiarism can lead to a reduced or fail
grade, in serious cases, for the Graduation Project course.

\vspace{0.8cm}
\noindent
\begin{tabular}{@{}p{0.48\textwidth}p{0.48\textwidth}@{}}
\textbf{Student name:} Omar Hossam Eldin Mohamed Hafez & \textbf{Student name:} Mohamed Sameh Mohamed Mohamed El Bassuni\\
Registration number: 20101737 & Registration number: 18100896\\
Signed: \rule{3cm}{0.3pt} & Signed: \rule{3cm}{0.3pt}\\
Date: \rule{3cm}{0.3pt} & Date: \rule{3cm}{0.3pt}\\[0.6cm]
\textbf{Student name:} Mohamed Mohamed Talaat Abokhatwa & \textbf{Student name:} Ahmed Abdelaal Saad\\
Registration number: 12103172 & Registration number: 16101027\\
Signed: \rule{3cm}{0.3pt} & Signed: \rule{3cm}{0.3pt}\\
Date: \rule{3cm}{0.3pt} & Date: \rule{3cm}{0.3pt}\\[0.6cm]
\textbf{Student name:} Adam Mohamed Ibrahim & \\
Registration number: 20101978 & \\
Signed: \rule{3cm}{0.3pt} & \\
Date: \rule{3cm}{0.3pt} & \\
\end{tabular}

% ─────────────────────────── PROJECT SUMMARY ───────────────────────────
\clearpage
\section*{Senior Project Summary Report}
\addcontentsline{toc}{section}{Senior Project Summary Report}

\begin{center}
Arab Academy for Science, Technology and Maritime Transport\\
College of Engineering and Technology --- Department of Computer Engineering\\
Academic Year: 2025/2026 \qquad Semester: 10
\end{center}

\renewcommand{\arraystretch}{1.3}
\begin{longtable}{@{}p{0.24\textwidth}p{0.72\textwidth}@{}}
\toprule
\textbf{Project title} & \projecttitle\\
\midrule
\textbf{Supervisor(s)} & Yahia Elderiny; Rania Kadry\\
\midrule
\textbf{Team members} &
Omar Hossam Eldin Mohamed Hafez --- 20101737\newline
Mohamed Sameh Mohamed Mohamed El Bassuni --- 18100896\newline
Mohamed Mohamed Talaat Abokhatwa --- 12103172\newline
Ahmed Abdelaal Saad --- 16101027\newline
Adam Mohamed Ibrahim --- 20101978\\
\midrule
\textbf{Project deliverables} &
A physics-constrained fault injection engine for OBD-II time series; a
six-class XGBoost fault classifier with SHAP explanations; a PID forecaster
and one-class anomaly detector; an Android data-collection app (OBD-LOG); a
Streamlit dashboard with an end-of-read health report; and a documented
protocol for future real-fault data collection.\\
\midrule
\textbf{Description} &
The system reads live engine sensor data through the standard OBD-II port at
1\,Hz, converts 60-second windows into 83 vehicle-agnostic features, and
classifies the engine state into healthy, cold-start, or one of four
developing fault classes (vacuum leak, fuel delivery, coolant sensor,
throttle sensor). Because no public real-fault OBD-II dataset exists at this
sample rate, training data is produced by injecting physically plausible
fault signatures into 84.5 hours of real healthy driving data. The pipeline
transfers across vehicles by z-scoring every feature against the target
car's own healthy baseline, and it reports ``Untested'' for any fault whose
required sensor a given car does not expose.\\
\midrule
\textbf{Engineering standards} &
SAE J1979 (OBD-II diagnostic services and PID definitions); ISO 15765-4
(CAN transport for diagnostics); ELM327 AT command set for adapter control;
OBD-II physical ranges used as hard clamps in the injection engine (fuel
trim $\pm$25\,\%, timing advance $\pm$64$^{\circ}$).\\
\midrule
\textbf{Design constraints} &
1\,Hz sampling ceiling of consumer ELM327 adapters (measured 0.37--1.10\,Hz
depending on PID count); Bluetooth Classic SPP throughput; per-vehicle PID
availability (the test Yaris exposes no MAP sensor); no labeled real-fault
data available anywhere at 1\,Hz; laptop-only inference (no cloud
dependency for the demo); all injected values clamped to engine physics.\\
\midrule
\textbf{Project impact} &
Earlier warning of developing engine faults than the check-engine light,
using hardware the average driver can afford (a phone plus a
\$10 adapter); a reusable, physics-grounded methodology for synthesizing
labeled fault data where none exists; an honest evaluation framework that
reports what it cannot test instead of fabricating scores.\\
\midrule
\textbf{Team organization} &
Omar Hossam Eldin Mohamed Hafez --- data pipeline and fault injection.\newline
Mohamed Sameh Mohamed Mohamed El Bassuni --- Android application and vehicle testing.\newline
Mohamed Mohamed Talaat Abokhatwa --- feature engineering and classifier.\newline
Ahmed Abdelaal Saad --- forecaster, anomaly detection, evaluation.\newline
Adam Mohamed Ibrahim --- dashboard, integration, documentation.\\
\midrule
\textbf{Ethics / safety} &
No physical fault was induced on any test vehicle. All driving data was
collected on public roads with the driver's consent, at legal speeds, with
the phone mounted and operated by a passenger. Model outputs are advisory
only; a qualified mechanic must confirm any suspected fault before repair
decisions. The synthetic origin of all fault labels is disclosed in every
result table.\\
\bottomrule
\end{longtable}

\vspace{0.6cm}
\noindent Main Supervisor Signature: \rule{5cm}{0.3pt}

% ─────────────────────────── STUDENTS CONTRIBUTION ───────────────────────────
\clearpage
\section*{Students Contribution}
\addcontentsline{toc}{section}{Students Contribution}

\begin{center}
\projecttitle\\[0.4cm]
Submitted by: Omar Hossam Eldin Mohamed Hafez, Mohamed Sameh Mohamed Mohamed El Bassuni, Mohamed Mohamed Talaat Abokhatwa, Ahmed Abdelaal Saad, Adam Mohamed Ibrahim
\end{center}

\vspace{0.5cm}
\begin{center}
\begin{tabular}{clp{5.5cm}}
\toprule
\textbf{Chapter} & \textbf{Title} & \textbf{Contributors}\\
\midrule
1 & Introduction & Omar Hossam Eldin Mohamed Hafez\newline Adam Mohamed Ibrahim\\
\addlinespace
2 & Literature Survey and Related Work & Mohamed Sameh Mohamed Mohamed El Bassuni\newline Mohamed Mohamed Talaat Abokhatwa\\
\addlinespace
3 & The Proposed System & Omar Hossam Eldin Mohamed Hafez\newline Mohamed Sameh Mohamed Mohamed El Bassuni\newline Mohamed Mohamed Talaat Abokhatwa\\
\addlinespace
4 & Experimental Results and Evaluation & Ahmed Abdelaal Saad\newline Adam Mohamed Ibrahim\\
\addlinespace
5 & Conclusion and Future Work & Ahmed Abdelaal Saad\newline Adam Mohamed Ibrahim\\
\bottomrule
\end{tabular}
\end{center}

% ─────────────────────────── ACKNOWLEDGMENT ───────────────────────────
\clearpage
\section*{Acknowledgment}
\addcontentsline{toc}{section}{Acknowledgment}

We wish to express our sincere gratitude to all those who gave us the
opportunity to make tangible progress in this project.

First and foremost, we thank God Almighty for giving us the strength,
knowledge, ability and opportunity to undertake this work and to persevere
and complete it satisfactorily.

Special appreciation for the supervision, advice and continuous encouragement
goes to our advisors, Yahia Elderiny and Rania Kadry.

We would also like to thank the Arab Academy for Science, Technology and
Maritime Transport for the facilities and support that made this project
possible, and Eron J.\ Maranh\~ao for publishing the carOBD dataset that our
training pipeline is built on.

Last but not least, our heartfelt gratitude goes to our families for their
support throughout the year, and to the long-suffering 2015 Toyota Yaris that
idled through every data-collection session without complaint.

% ─────────────────────────── ABSTRACT ───────────────────────────
\clearpage
\section*{Abstract}
\addcontentsline{toc}{section}{Abstract}

Modern cars continuously report engine sensor data through the standardized
OBD-II port, yet that data is used almost exclusively after a fault has
already tripped a diagnostic trouble code. This project builds a predictive
maintenance system that watches the sensor stream itself and flags a
developing fault before the check-engine light fires. The system reads
14 engine parameters at 1\,Hz, summarizes each 60-second window into 83
vehicle-agnostic features, and classifies the window into healthy,
cold-start, or one of four fault classes: intake vacuum leak, fuel delivery
fault, coolant temperature sensor fault, and throttle position sensor fault.

No public OBD-II dataset with labeled real faults exists at this sample rate,
so we trained on synthetic faults injected into 84.5 hours of real healthy
driving data from the public carOBD corpus --- 129 sessions that we recovered
from an upstream CSV parsing defect which had previously made 120 of them
appear corrupt. The injection engine is physics-constrained: every written
value respects engine causality, from the 1$^{\circ}$C-per-second thermal
inertia of coolant to the short-term-to-long-term fuel trim handoff of a real
adaptive ECU. An XGBoost classifier reaches a leave-one-family-out macro-F1
of 0.977 (worst family 0.958), and we report these numbers explicitly as
synthetic self-consistency scores, not as proof of real-fault detection. On a
real, healthy 2015 Toyota Yaris the deployed pipeline correctly declared the
air-system class Untested (the car has no MAP sensor) and refused to issue a
verdict on insufficient evaluable data --- behavior we argue is an
engineering contribution in itself, because an earlier version of the same
pipeline confidently reported a 97\,\% air-system fault on the same healthy
car. The system ships with an Android data-collection app, a live Streamlit
dashboard, and a documented protocol for the real-fault recordings that would
convert the synthetic validation into field validation.

% ─────────────────────────── TOC / LOF / LOT ───────────────────────────
\clearpage
\tableofcontents
\clearpage
\listoffigures
\addcontentsline{toc}{section}{List of Figures}
\listoftables
\addcontentsline{toc}{section}{List of Tables}

% ─────────────────────────── ACRONYMS ───────────────────────────
\clearpage
\section*{List of Acronyms}
\addcontentsline{toc}{section}{List of Acronyms}

\begin{acronym}[XGBoost]
  \acro{OBD-II}{On-Board Diagnostics, second generation --- the standardized
    vehicle diagnostic interface mandated on passenger cars since 1996}
  \acro{PID}{Parameter ID --- one addressable sensor or computed value in the
    OBD-II protocol (e.g.\ 0x0C = engine RPM)}
  \acro{ELM327}{The de-facto standard microcontroller that bridges the OBD-II
    bus to a serial (Bluetooth/USB) link using AT commands}
  \acro{ECU}{Engine Control Unit --- the embedded computer that manages
    fueling, ignition and emissions}
  \acro{MAF}{Mass Air Flow sensor --- measures intake airflow directly}
  \acro{MAP}{Manifold Absolute Pressure sensor --- measures intake manifold
    pressure; the airflow reference on speed-density engines}
  \acro{XGBoost}{eXtreme Gradient Boosting --- a gradient-boosted decision
    tree library}
  \acro{SHAP}{SHapley Additive exPlanations --- a game-theoretic method for
    attributing a model prediction to its input features}
  \acro{ML}{Machine Learning}
  \acro{LOFO}{Leave-One-Family-Out cross-validation (holding out an entire
    session family, e.g.\ all idle recordings)}
  \acro{LOSO}{Leave-One-Session-Out cross-validation}
  \acro{LTFT}{Long-Term Fuel Trim --- the ECU's learned, persistent fuel
    correction, in percent}
  \acro{STFT}{Short-Term Fuel Trim --- the ECU's fast, oscillating fuel
    correction, in percent}
  \acro{DTC}{Diagnostic Trouble Code --- the standardized fault code that
    turns on the check-engine light}
  \acro{API}{Application Programming Interface}
  \acro{CSV}{Comma-Separated Values file format}
  \acro{RPM}{Revolutions Per Minute (engine speed)}
  \acro{IAT}{Intake Air Temperature}
  \acro{TPS}{Throttle Position Sensor}
  \acro{APK}{Android Package --- the installable Android application format}
  \acro{BOM}{Byte Order Mark --- a Unicode signature at the start of a text file}
  \acro{Kafka}{Apache Kafka --- a distributed event-streaming platform}
  \acro{REST}{REpresentational State Transfer --- an HTTP API architectural style}
\end{acronym}
\clearpage
